\chapter{Desarrollo de la comparativa}
\label{cap:comparativa}

% =====================================================================================
% GUION DEL CAPITULO — ESTRUCTURA FIJADA POR EL AUTOR (sep-2026).
%
% 🔴 NO AÑADIR SECCIONES SIN CONSULTARLO. La estructura acordada es exactamente esta:
%      7.1 Estrategia de entrenamiento   7.3 Random Forest
%      7.2 Ridge                         7.4 GNN
%
% ⚠️ La version anterior del capitulo se DESCARTO entera (sep-2026): duplicaba el
%    capitulo~\ref{cap:planteamiento} —tenia sus propias secciones de particion,
%    preprocesado y metricas— y ademas repetia la explicacion de los tres modelos. Si
%    hace falta recuperar algo, esta en el historial del repositorio.
%
% 🔴 REPARTO CON EL CAPITULO 6, que es la regla que evita volver a duplicar:
%      cap. 6 -> los DATOS y el diseño experimental: de donde salen, que dice la EDA,
%                como se preprocesan y con que metricas se juzgan.
%      cap. 7 -> los MODELOS: cual es cada uno, por que esta, con que hiperparametros y
%                como se ha entrenado.
%    Si una frase de aqui habla del conjunto de datos, del umbral o de una metrica, esta
%    en el capitulo equivocado: se cita el apartado del 6 y se sigue.
%
% ⚠️ EXTENSION: el capitulo NO debe ser largo. Una figura por modelo, y la tabla de las
%    tres convoluciones al final de §7.4.
%
% FUENTES, por seccion:
%   7.1 -> notebooks/03_ridge_rf/NOTAS_ENTRENAMIENTO.md
%          notebooks/04_gnn/NOTAS_ENTRENAMIENTO.md
%          doc/changes/2026-08_cambio_decisiones.md §4 (la escalera de seis peldaños)
%   7.2 -> src/baselines.py · scripts/train_ridge.py
%   7.3 -> src/baselines.py · scripts/train_rf.py · liao2024machine (el mejor tabular)
%   7.4 -> src/gnn_model.py · src/train.py · configs/gnn_config.yaml
%          notebooks/04_gnn/NOTAS_ENTRENAMIENTO.md · doc/changes/2026-08_cambio_decisiones.md §9
%
% 🔴 CIFRAS PROHIBIDAS (son del dataset v1, obsoleto): «30-70 saltos» de distancia entre
%    nodos —la del v2 es 44— y «42,7x» / «45,4x» del QFT —la citable es 3,69x—.
%
% ⚠️ NINGUN RESULTADO ENTRA AQUI. Este capitulo dice como se entrena; que salio de ello
%    es el capitulo~\ref{cap:resultados}. Si aparece un R² o un MAE, esta mal colocado.
% =====================================================================================


% =====================================================================================
\section{Estrategia de entrenamiento}
\label{sec:entrenamiento}
% =====================================================================================

La estrategia del entrenamiento se compone de estos seis peldaños:

\begin{enumerate}
    \item \textbf{No mitigar.} El suelo: si no se bate, el método empeora las cosas y no sirve.
    \item \textbf{Factor de supervivencia constante}, la media de entrenamiento. ¿Aporta saber la media?
    \item \textbf{Ridge con una sola variable}, la duración total. ¿Aporta saber cuánto dura
          el circuito?
    \item \textbf{Ridge con todas las variables.} ¿Aporta solo la linealidad?
    \item \textbf{Random Forest.} ¿Aporta la no linealidad?
    \item \textbf{Red de grafos.} ¿Aporta la estructura del circuito?
\end{enumerate}

Los tres modelos que se comparan son los tres últimos peldaños y cada uno añade una capacidad sobre el anterior: Ridge es el
\textbf{baseline} y fija el suelo, Random Forest añade no \textbf{linealidad} y está porque es el mejor modelo tabular de los que compara la referencia \cite{liao2024machine}. La red de grafos cambia la representación de los datos, y la diferencia entre ella y los dos primeros es \textbf{el valor de la
estructura que la agregación descarta}.

Cabe recalcar una asimetría presente. Ridge y la red de
grafos se entrenan con una \textbf{pérdida de Huber}, que penaliza los errores grandes de forma lineal y evita que los factores de supervivencia extremos se sesgen el entrenamietno. Sin embargo, \textbf{Random Forest no la admite} y no se emplea.


% =====================================================================================
\section{Ridge}
\label{sec:sol_ridge}
% =====================================================================================

Una regresión lineal con penalización $L_2$ \cite{hoerl1970ridge}. No se usan mínimos cuadrados
básico por un motivo medido en el apartado~\ref{sec:colinealidad}: aun después de descartar las
veinte columnas redundantes y se prefiere que el modelo recaba información de ella ajustándose solo.

Se ajusta \textbf{un modelo por observable}, cinco en total, sobre las $34$ columnas
preprocesadas y descartando las filas sin etiqueta de es observable. El único
hiperparámetro es la constante de regularización, barrida en diez valores de $10^{-3}$ a
$10^{6}$. El óptimo cae en $100$ para la media de $Z$ y el correlador de vecinos, y en $1{.}000$
para los otros tres.

La Figura~\ref{fig:ridge_coef} recoge los doce coeficientes de mayor magnitud. Se muestran
\textbf{con su signo}, que es lo único que este modelo da y los otros dos no: la dirección en
la que cada variable empuja al factor de supervivencia.

% ⚠️ FORMATO: titulo arriba y fuente abajo (Instrucciones §1.3, p. 6).
% Figura propia. La genera modelos.dibujar_ridge_coeficientes() y se copia desde
% TFM-Quantum/figures/modelos/ridge_coeficientes.pdf. Sin titulo interno.
\begin{figure}[htbp]
    \centering
    \caption{Los doce coeficientes mayores de Ridge, con su signo, promediando los cinco
             observables.}
    \label{fig:ridge_coef}

    \includegraphics[width=0.92\textwidth]{cap07/ridge_coeficientes}

    {\footnotesize Fuente: elaboración propia.}
\end{figure}

% =====================================================================================
\section{Random Forest}
\label{sec:sol_rf}
% =====================================================================================

Un bosque de árboles de decisión \cite{breiman2001random} añade sobre Ridge no linealidad.

Cabe recalcar que este modelo \textbf{no extrapola} y no será util para comparar el \textbf{eje tamaño}.

Se entrena igual que Ridge (un modelo por observable, las mismas $34$ columnas y los mismos
descartes). El escalado no le afecta, pero \textbf{se le aplican las mismas transformaciones}: no por necesidad, sino para que la única diferencia entre los dos
modelos tabulares sea el modelo. La rejilla fija $300$ árboles y barre la profundidad máxima
($\{$sin límite, $6$, $12$, $24\}$), las hojas mínimas ($\{1, 5, 20, 50\}$) y las variables por
división ($\{\sqrt{\cdot}, 0{,}5, 1{,}0\}$). Los óptimos salen \textbf{poco profundos}: profundidad $6$ en tres
de los cinco observables, con hojas de $5$ o $50$ muestras y la mitad de las variables por
división. En la Figura~\ref{fig:rf_imp} se muestra que el modelo se apoya en muy pocas columnas.

% Figura propia. La genera modelos.dibujar_rf_importancias() y se copia desde
% TFM-Quantum/figures/modelos/rf_importancias.pdf. Sin titulo interno.
\begin{figure}[htbp]
    \centering
    \caption{Las doce variables más importantes para Random Forest, promediando los cinco
             observables.}
    \label{fig:rf_imp}

    \includegraphics[width=0.92\textwidth]{cap07/rf_importancias}

    {\footnotesize Fuente: elaboración propia.}
\end{figure}

Las dos figuras cuentan lo mismo y coinciden con que \textbf{la duración total y la profundidad tiene mayor peso en el modelo},
con un $0{,}39$ y un $0{,}33$ de importancia frente al $0{,}03$ del siguiente.


% =====================================================================================
\section{La red de grafos}
\label{sec:gnn}
% =====================================================================================

Una red de paso de mensajes sobre el grafo del circuito. Es la que más justificación necesita,
porque el diseño se apartó de la arquitectura de referencia del área.

\textbf{Paso de mensajes y no atención.} El trabajo más cercano \cite{bao2025gtraqem} usa un
\textit{graph transformer} \emph{sin} paso de mensajes; sin embargo, el computo disponible era insuficiente y por eso ser recurrió a modelos de grafos menos pesados.

\textbf{El resumen global es la pieza que hace funcionar el modelo}
(Figura~\ref{fig:gnn_arq}). Se escogió incorporar un vector que conecte a todos lo nodos para ahorra la distancia de paso entre cualquier nodo y, así, usar ese vector como atajo. La idea se toma del mismo trabajo\cite{bao2025gtraqem}, con un cambio: allí el
vector es \textbf{unidireccional}, y eso basta en un \textit{transformer} porque su atención ya
conecta todo con todo; sin embargo, en una red de paso de mensajes no, y lo que queda es un \textit{pooling} aprendido calculando con
\textbf{media}. Los grafos entran \textbf{completos}.

% FIGURA PROPIA en TikZ. Codigo en figuras/cap07/gnn_arquitectura_tikz.tex.
% ⚠️ Lo que la figura tiene que dejar claro son las DOS flechas del resumen global: lee y
%    devuelve. Es lo que lo distingue de un pooling aprendido, y en prosa no se ve.
\begin{figure}[!ht]
    \centering
    \caption{La arquitectura. El resumen global lee de todos los nodos y les devuelve: son
             las dos flechas, no una.}
    \label{fig:gnn_arq}

    % =====================================================================================
% La arquitectura del MPNN — figura propia en TikZ.
%
% Lo que tiene que quedar claro y no se entiende sin dibujo: el resumen global es
% BIDIRECCIONAL. Lee de todos los nodos y DEVUELVE a todos los nodos, y esa flecha de
% vuelta es la que lo distingue de un pooling aprendido. Por eso hay dos flechas entre la
% capa de mensajes y el resumen, y no una.
%
% CODIGO DE COLOR, el de la memoria:
%   azulunir  -> el dato (el grafo de entrada y la salida)
%   acentofig -> lo que se le hace (las capas y el resumen)
%
% ⚠️ \shorthandoff{<>} es obligatorio: `babel` en español hace activos < y >.
% =====================================================================================
\begingroup
\shorthandoff{<>}%
\hyphenpenalty=10000\exhyphenpenalty=10000\relax

\begin{tikzpicture}[
    >={Stealth[length=2mm]},
    font=\footnotesize,
    dato/.style  = {draw=azulunir, line width=0.9pt, rounded corners=3pt, fill=azulunir!6,
                    align=center, inner sep=5pt, text width=2.1cm},
    capa/.style  = {draw=acentofig, line width=0.9pt, rounded corners=3pt,
                    fill=acentofig!8, align=center, inner sep=5pt, text width=2.3cm},
    resumen/.style = {draw=acentofig, line width=1.2pt, rounded corners=3pt,
                      fill=acentofig!16, align=center, inner sep=5pt, text width=2.6cm},
  ]

  % --- la cadena principal -------------------------------------------------------------
  \node[dato]  (entrada) at (0,0)
    {\textbf{Grafo}\\{\scriptsize $21$ dims por nodo}};
  \node[capa]  (mp)      at (3.4,0)
    {\textbf{Paso de mensajes}\\{\scriptsize $4$ capas · aristas dirigidas}};
  \node[capa]  (lectura) at (7.2,0)
    {\textbf{Lectura}\\{\scriptsize resumen $\oplus$ \textit{pooling}}};
  \node[dato]  (salida)  at (10.4,0)
    {\textbf{$\hat f$}\\{\scriptsize cinco salidas}};

  \draw[->, azulunir,  line width=1pt] (entrada.east) -- (mp.west);
  \draw[->, acentofig, line width=1pt] (mp.east)      -- (lectura.west);
  \draw[->, azulunir,  line width=1pt] (lectura.east) -- (salida.west);

  % --- el resumen global, encima, con SUS DOS flechas ------------------------------------
  \node[resumen] (virtual) at (3.4,2.0)
    {\textbf{Resumen global}\\{\scriptsize un vector por grafo · pesos propios}};

  \draw[->, acentofig, line width=1pt] ($(mp.north)+(-0.42,0)$) -- ($(virtual.south)+(-0.42,0)$)
        node[midway, left=1pt, font=\scriptsize, text=black] {lee};
  \draw[->, acentofig, line width=1pt] ($(virtual.south)+(0.42,0)$) -- ($(mp.north)+(0.42,0)$)
        node[midway, right=1pt, font=\scriptsize, text=black] {devuelve};

  \draw[->, acentofig, line width=0.8pt, dotted]
        (virtual.east) to[out=0, in=110] (lectura.north);

  % --- la nota que justifica la vuelta ---------------------------------------------------
  \node[align=left, font=\scriptsize, text=gray, text width=3.7cm] at (8.9,2.1)
    {con $4$ capas un nodo alcanza\\el $1{,}25$\,\% de los pares:\\el resumen es la única vía\\al largo alcance};

  % --- lo que envuelve cada capa ---------------------------------------------------------
  \node[align=center, font=\scriptsize, text=gray] at (3.4,-1.35)
    {cada capa: normalización + conexión residual};

\end{tikzpicture}%
\endgroup


    {\footnotesize Fuente: elaboración propia.}
\end{figure}

\textbf{Tres convoluciones, en orden de complejidad creciente} (Tabla~\ref{tab:convs}): se sube
de peldaño solo si el anterior se queda corto, de modo que si la atención local mejora se sabe
que aporta, y si no, también.

\begin{table}[!ht]
    \centering
    \caption{Las tres convoluciones que se prueban y la configuración con la que se entrenó
             cada una.}
    \label{tab:convs}
    {\footnotesize
    \begin{tabular}{@{}p{0.20\textwidth}p{0.27\textwidth}rrr@{}}
        \toprule
        \textbf{Convolución} & \textbf{Qué añade} & \textbf{Oculta} & \textbf{Tasa}
            & \textbf{Parámetros} \\
        \midrule
        \texttt{SAGEConv} \cite{hamilton2017sage}
            & Agrega los vecinos por media. Ligera y estable: es la de arranque
            & $256$ & $10^{-3}$ & $1{,}06$~M \\
        \texttt{GINConv} \cite{xu2019gin}
            & Suma estricta más un perceptrón interno. La más expresiva; distingue mejor
              topologías parecidas
            & $256$ & $3\cdot10^{-4}$ & $1{,}06$~M \\
        \texttt{GATv2Conv} \cite{brody2022gatv2}
            & \textbf{Atención local}: aprende a pesar más a los vecinos con peor telemetría
            & $160$ & $10^{-3}$ & $1{,}04$~M \\
        \bottomrule
    \end{tabular}}

    {\footnotesize Fuente: elaboración propia. Las tres con cuatro capas.}
\end{table}

\textbf{Las tres variante se prueban con el mismo presupuesto de
parámetros}

El resto de la rejilla barre la tasa de aprendizaje ($\{3\cdot10^{-4}, 10^{-3},
3\cdot10^{-3}\}$), la dimensión oculta ($\{128, 192, 256\}$) y el número de capas ($\{4, 6\}$),
con lotes de $32$ grafos y regularización de $10^{-5}$. El entrenamiento se hace en una máquina
con tarjeta gráfica en la nube \cite{google2026cloud}, porque en local no hay ninguna.